\chapter{Fields}

\cmd{FIELD} defines generic tabular data. Fields may be defined at root level or inside a loadcase block, depending on their intended use. The field domain determines how rows are associated with model entities.

\section{Supported domains}

For node, element, integration-point, and material-point fields the general form is

\begin{syntax}
*FIELD, NAME=<name>, TYPE=NODE|ELEMENT|ELEMENT_IP|ELEMENT_MP, COLS=<n>, FILL=ZERO|NAN
<row-id>, <v1>, <v2>, ...
\end{syntax}

Element-nodal fields use an element identifier and zero-based local node index instead of one flat row identifier:

\begin{syntax}
*FIELD, NAME=<name>, TYPE=ELEMENT_NODAL, COLS=<n>, FILL=ZERO|NAN
<element-id>, <local-node>, <v1>, <v2>, ...
\end{syntax}

\key{COLS} is required and must be greater than zero. FEMaster limits fields to 64 components. \key{FILL} defaults to \val{ZERO}; \val{NAN} initializes unspecified values with NaN. Individual values may be omitted; omitted components preserve the initialization value already stored in the field.

\begin{longtable}{@{}p{0.20\linewidth}p{0.70\linewidth}@{}}
\toprule
\textbf{Domain} & \textbf{Meaning} \\
\midrule
\endhead
\val{NODE} & One row per node ID. Used for temperature fields, initial velocities, or other prescribed nodal data. \\
\val{ELEMENT} & One row per element ID. Used for topology density and orientation fields. \\
\val{ELEMENT_NODAL} & One row per element-local node. Input rows are addressed by element ID and zero-based local node index. Used, for example, for prescribed shell director fields. \\
\val{ELEMENT_IP} & One row per enumerated element integration point. The alias \val{IP} is accepted. \\
\val{ELEMENT_MP} & One row per enumerated material point. The alias \val{MP} is accepted. This domain is also used by nonlinear material-state infrastructure. \\
\bottomrule
\end{longtable}

The long aliases \val{ELEMENTNODAL}, \val{ELEMENTIP}, and \val{ELEMENTMP} are accepted by the domain parser; the underscore-separated forms are recommended in input decks.

Enumeration-dependent domains require the corresponding element topology and local enumeration to exist. FEMaster therefore creates these fields after topology construction during its field parsing stage.

Data lines may omit individual values; omitted values do not overwrite existing entries. The tokens \val{NAN}, \val{+NAN}, \val{-NAN}, \val{INF}, \val{+INF}, and \val{-INF} are accepted.

\section{Typical references}

\begin{longtable}{@{}p{0.25\linewidth}p{0.25\linewidth}p{0.4\linewidth}@{}}
\toprule
\textbf{Use} & \textbf{Domain/Cols} & \textbf{Referencing command} \\
\midrule
\endhead
Temperature load & \val{NODE}, 1 & \cmd{TLOAD} with \key{TEMPERATUREFIELD}. \\
Initial velocity & \val{NODE}, 6 & \cmd{INITIALVELOCITY}. \\
Topology density & \val{ELEMENT}, 1 & \cmd{TOPODENSITY} with \key{FIELD=<name>}. \\
Topology orientation & \val{ELEMENT}, 3 & \cmd{TOPOORIENT} with \key{FIELD=<name>}. \\
Shell directors & \val{ELEMENT_NODAL}, 3 & \cmd{NORMAL} with \key{FIELD=<name>}. \\
\bottomrule
\end{longtable}

\section{Shell normal selection: \cmd{NORMAL}}

\begin{syntax}
*NORMAL, FIELD=<element-nodal-field>
\end{syntax}

\cmd{NORMAL} selects an existing three-component \val{ELEMENT_NODAL} field as the shell-normal field. It is valid at model root scope. The command does not allocate the field and does not create, normalize, or average the supplied vectors; geometric completion remains part of shell-normal construction.

A typical partially prescribed field uses \key{FILL=NAN} so unspecified element-local directions remain distinguishable from explicitly supplied vectors:

\begin{exampledeck}
*FIELD, NAME=DIRECTORS, TYPE=ELEMENT_NODAL, COLS=3, FILL=NAN
10, 0, 0.0, 0.0, 1.0
10, 1, 0.0, 0.0, 1.0
*NORMAL, FIELD=DIRECTORS
\end{exampledeck}

The \cmd{FIELD} definition must appear before the referencing \cmd{NORMAL} command. Further shell-specific interpretation is described in Section~\ref{sec:shell-directors}.
